\E{ignidrome}{ˈɪɡnɪdrəʊm}{n.}{L. \la{ignis} fire + Gk. δρόμος course} The path of a flame through fuel; \lbl{fig.} the course of an innovation through a society, which burns fastest where the material is driest and the regulation thinnest. \ex{Trace the ignidrome and you will find the frontier.} \src{Hom. Grad.} iii. --- \dv{ignidromic} \textit{adj.}

\E{illuminarch}{ɪˈluːmɪnɑːk}{n.}{L. \la{illuminare} to light up + Gk. ἀρχός ruler} One who governs by making things visible; a publisher of data, a builder of instruments, an opener of archives. Opp. \xr{katechontist}. --- \dv{illuminarchy} \textit{n.}

\E{immoratory}{ɪˈmɒrətəri}{adj.}{L. \la{in-} not + \la{mora} delay} Admitting of no delay; said of decisions, permits and launches that are granted on the day they are asked. \ex{an immoratory licence, rarer than a comet}. --- \dv{immoratorily} \textit{adv.} Opp. \xr{moratorian}.

\E{impedimentarian}{ɪmˌpɛdɪmɛnˈtɛəriən}{n.}{L. \la{impedimentum} hindrance, baggage + -arian} \sn{1} One whose profession is the raising of obstacles. \sn{2} One who travels with so much procedural baggage that the journey is never begun. \ex{The impedimentarians arrived with forty trunks of process and no destination.} \src{Sat. Decel.} x. --- \dv{impedimentarianism} \textit{n.}

\E{inaugurist}{ɪˈnɔːɡjʊrɪst}{n.}{L. \la{inaugurare} to take omens, to consecrate a beginning} One who delights in beginnings; a serial founder, more drawn to the first day of a venture than to its thousandth. --- \dv{inaugurism} \textit{n.} Cf. \xr{neoctist}.

\E{incubatrice}{ˈɪŋkjʊbətriːs}{n.}{Fr. \la{incubatrice}, f. L. \la{incubare} to lie upon, hatch} An institution for the hatching of ventures; an accelerator of companies, figured as a warm mother-machine. \ex{Every incubatrice is a small wager against the bradyerie.} \src{Ep. Fabr.} v. --- \dv{incubatricial} \textit{adj.}

\E{infinitivore}{ɪnˈfɪnɪtɪvɔː}{n.}{L. \la{infinitus} unbounded + \la{vorare} to devour} A devourer of the infinite; one whose appetite for progress no achievement can satisfy. --- \dv{infinitivorous} \textit{adj.} Cf. \xr{apeirophagy}.

\E{iterophily}{ˌɪtəˈrɒfɪli}{n.}{L. \la{iterum} again + Gk. φιλία love} The love of iteration; the conviction that the hundredth draft, shipped, is worth more than the perfect plan, unwritten. \ex{Iterophily is humility made fast.} \src{Ep. Fabr.} iii. 2. --- \dv{iterophile} \textit{n.} Cf. \xr{ecpempsis}.

\E{jaculatrix}{ˌdʒækjʊˈleɪtrɪks}{n., pl. -trices}{L. \la{jaculari} to hurl + fem. agentive -trix} A launching-engine; a rocket, catapult or mass-driver, considered as a thrower of things toward the sky. --- \dv{jaculatory} \textit{adj.} pertaining to launch. Cf. \xr{ouranobasis}.

\E{juvenocracy}{ˌdʒuːvɪˈnɒkrəsi}{n.}{L. \la{juvenis} young + Gk. κράτος rule} Government by the young; the polity of the garage and the dormitory, in which credentials are replaced by commits. Opp. \xr{gerontostasis}. --- \dv{juvenocrat} \textit{n.}

\E{kalotechny}{ˈkæləʊtɛkni}{n.}{Gk. καλός beautiful + τέχνη craft} The beauty of well-made machines; the aesthetic which finds in a turbine or a clean codebase the dignity once reserved for temples. --- \dv{kalotechnic} \textit{adj.}

\E{kinesiarch}{kɪˈniːsiɑːk}{n.}{Gk. κίνησις motion + ἀρχός leader} A prime mover in a small way; the person in any organisation who causes things to happen. \ex{A firm without a kinesiarch is a museum with payroll.} \src{Ep. Fabr.} vii. Cf. \xr{impetuarch}.

\E{kosmergate}{ˈkɒzmɜːɡeɪt}{n.}{Gk. κόσμος world + ἐργάτης worker} A labourer on the cosmos; one engaged in works of planetary or stellar scale. \ex{The kosmergates do not ask whether the sky is the limit; they ask what it is made of.} \src{Lib. Ascens.} vi.

\E{laborifuge}{ləˈbɒrɪfjuːdʒ}{n.}{L. \la{labor} toil + \la{fugere} to flee} One who flees work; \lbl{spec.} one who prefers to write a report on a problem rather than to solve it. Opp. \xr{ponophily}.

\E{lampadarch}{ˈlæmpədɑːk}{n.}{Gk. λαμπαδάρχης leader of a torch-race} The leader of a relay of builders; one who receives a project from predecessors and hands it on, enlarged, to successors. Cf. \xr{lampadephory}.

\E{latentist}{ˈleɪtəntɪst}{n.}{L. \la{latens} lying hidden} One who holds that the latent capabilities of a model, a material or a people exceed by far what has been demonstrated, and who acts on that belief before the demonstration. --- \dv{latentism} \textit{n.} Cf. \xr{anaptyxis}.

\E{legiferomania}{lɛˌdʒɪfərəʊˈmeɪniə}{n.}{L. \la{legifer} law-bearing + Gk. μανία madness} The compulsive production of statutes; a fever of lawmaking which mistakes the length of the code for the strength of the state. \ex{In the year of the legiferomania they passed nine hundred acts and built one bridge.} \src{Disp. Mor.} ii. --- \dv{legiferomaniac} \textit{n.} Cf. \xr{regulomania}.

\E{leptotechny}{ˈlɛptəʊtɛkni}{n.}{Gk. λεπτός fine, slender + τέχνη craft} The craft of the very small; engineering at the scale of atoms, from the transistor to the molecular machine. --- \dv{leptotechnician} \textit{n.} Cf. \xr{lithurgist}.

\E{lethargocracy}{ˌlɛθəˈɡɒkrəsi}{n.}{Gk. λήθαργος drowsy + κράτος rule} Government by drowsiness; a polity that is not opposed to progress but cannot remain awake long enough to approve it. --- \dv{lethargocrat} \textit{n.} Cf. \xr{acediocracy}.

\E{ludergy}{ˈluːdɜːdʒi}{n.}{L. \la{ludus} play + Gk. ἔργον work} The energy of play; the motive power of the hacker, who builds because building is a game, and who has accomplished more than the dutiful. \ex{No grant was ever as productive as a weekend of ludergy.} \src{Chron. Fabr.}

\E{luminivore}{luːˈmɪnɪvɔː}{n.}{L. \la{lumen} light + \la{vorare} to devour} An eater of light; any photosynthetic or photovoltaic system, from the leaf to the orbital array. --- \dv{luminivorous} \textit{adj.} Cf. \xr{heliophage}.

\E{lysinomy}{laɪˈsɪnəmi}{n.}{Gk. λύσις a loosing + νόμος law} The repeal of laws; the art of loosening the legal knots tied by former generations. \ex{Lysinomy is the most neglected branch of legislation.} \src{Tract. Cop.} viii. --- \dv{lysinomist} \textit{n.} Opp. \xr{nomosclerosis}.

\E{mantoclast}{ˈmæntəʊklæst}{n.}{Gk. μάντις prophet + κλάστης breaker} A breaker of prophecies; one who wagers against the predictions of doom and collects. Opp. \xr{olethromancer}. --- \dv{mantoclasm} \textit{n.}

\E{mechanogony}{ˌmɛkəˈnɒɡəni}{n.}{Gk. μηχανή machine + -γονία begetting} The begetting of machines by machines; the stage of industry at which factories build factories and models train models. --- \dv{mechanogonic} \textit{adj.} Cf. \xr{autocatalyse}.

\E{megalurgy}{ˈmɛɡəlɜːdʒi}{n.}{Gk. μεγαλουργία the doing of great works} The undertaking of great works; the building of things too large to be finished within a single term of office. \ex{Our ancestors practised megalurgy; we practise consultation.} \src{Hom. Grad.} iv. --- \dv{megalurge} \textit{n.}; \dv{megalurgic} \textit{adj.}

\E{melioritude}{ˌmiːliˈɒrɪtjuːd}{n.}{L. \la{melior} better + -itude} The condition of steadily becoming better; the felt quality of a life lived in an improving age.

\E{mellontarch}{məˈlɒntɑːk}{n.}{Gk. τὸ μέλλον the future + ἀρχός ruler} One who governs in the name of the future, representing the interests of the unborn against the comfort of the living. --- \dv{mellontarchy} \textit{n.} Cf. \xr{mellontolatry}.

\E{mesostasis}{mɛˈsɒstəsɪs}{n.}{Gk. μέσος middle + στάσις standing still} A halt half-way up the curve; the plateau upon which a civilization pauses, mistaking a landing for a summit. \ex{The mesostasis is comfortable, and therein lies its danger.} \src{Lib. Ascens.} iv. --- \dv{mesostatic} \textit{adj.} Cf. \xr{asymptophobia}.

\E{metabathmic}{ˌmɛtəˈbæθmɪk}{adj.}{Gk. μετά between + βαθμός step} Between rungs; of a civilization in transition from one level of energy capture to the next, as ours from the fossil to the nuclear. See \xr{anabathmism}.

\E{metakinesis}{ˌmɛtəkɪˈniːsɪs}{n.}{Gk. μετά beyond + κίνησις motion} Change in the manner of change itself; a revolution in how revolutions happen. --- \dv{metakinetic} \textit{adj.}

\E{myriergy}{ˈmɪriɜːdʒi}{n.}{Gk. μυρίοι countless + ἔργον work} The countless small works of an open economy, which together outbuild any single plan. \ex{The planner sees one dam; the market sees a myriergy.} \src{Tract. Cop.} ix.

\E{mythotechny}{ˈmɪθəʊtɛkni}{n.}{Gk. μῦθος story + τέχνη craft} The craft of telling stories that build things; the manufacture of the futures that draw investment toward themselves. --- \dv{mythotechnist} \textit{n.} Cf. \xr{hyperstition}.

\E{naupegy}{nɔːˈpiːdʒi}{n.}{Gk. ναυπηγία ship-building} The building of ships; \lbl{spec.} of vessels for the void. \ex{The first starship will be launched from a yard whose naupegy began with a spreadsheet.} \src{Lib. Ascens.} viii. --- \dv{naupegist} \textit{n.}

\E{necrostasis}{nɛˈkrɒstəsɪs}{n.}{Gk. νεκρός dead + στάσις standing still} The deathly stillness of an institution that persists after its purpose has gone; an office that outlives its reason. --- \dv{necrostatic} \textit{adj.} Cf. \xr{sclerocrat}.

\E{nephelarch}{ˈnɛfəlɑːk}{n.}{Gk. νεφέλη cloud + ἀρχός ruler} A lord of the cloud; the proprietor of a great computing utility. --- \dv{nephelarchy} \textit{n.} Cf. \xr{logisterion}.

\E{neophage}{ˈniːəfeɪdʒ}{n.}{Gk. νέος new + -φάγος eating} One who devours novelties; an early adopter. \ex{The neophage installs before reading and reads before sleeping.} \src{Agora} --- \dv{neophagous} \textit{adj.}

\E{nesopolis}{nɪˈsɒpəlɪs}{n.}{Gk. νῆσος island + πόλις city} A chartered city set apart, like an island, from the laws around it, for the trial of new institutions. --- \dv{nesopolitan} \textit{adj.} Cf. \xr{nautopolis}.

\E{noëtergy}{nəʊˈɛtɜːdʒi}{n.}{Gk. νοητός intelligible + ἔργον work} The energy expended in thought; the joules per inference. \ex{Noëtergy is the price of mind, and it is falling.} \src{Brev. Therm.} v.

\E{nomoclast}{ˈnɒməklæst}{n.}{Gk. νόμος law + κλάστης breaker} A breaker of laws for the sake of progress; one who builds first and litigates later. --- \dv{nomoclasm} \textit{n.} Cf. \xr{antiphragmatist}.

\E{nyctergy}{ˈnɪktɜːdʒi}{n.}{Gk. νύξ, νυκτός night + ἔργον work} Work done by night; the labour of builders across time-zones and of machines that never sleep. --- \dv{nyctergic} \textit{adj.} Cf. \xr{agrypnist}.

\E{ochlotachy}{ˈɒkləʊˌtæki}{n.}{Gk. ὄχλος crowd + τάχος speed} The speed of the crowd; the velocity with which an idea, true or false, spreads upon the Agora. --- \dv{ochlotachic} \textit{adj.} Cf. \xr{doxomachy}.

\E{oknocracy}{ɒkˈnɒkrəsi}{n.}{Gk. ὄκνος hesitation + κράτος rule} Government by hesitation; rule of the long pause before any decision. \ex{In the oknocracy, the question is always asked and never answered.} \src{Sat. Decel.} vii. --- \dv{oknocrat} \textit{n.}

\E{oneiroergy}{ˌəʊnaɪrəʊˈɜːdʒi}{n.}{Gk. ὄνειρος dream + ἔργον work} The conversion of dreams into works; the founder's chief and least documented skill.

\E{ophelocracy}{ˌɒfɪˈlɒkrəsi}{n.}{Gk. ὄφελος benefit, use + κράτος rule} Government by usefulness; a polity in which institutions are judged by what they deliver. --- \dv{ophelocrat} \textit{n.}

\E{orexergy}{ɒˈrɛksɜːdʒi}{n.}{Gk. ὄρεξις appetite + ἔργον work} The appetite for energy; the measure of a society's willingness to consume power in pursuit of its goals. Cf. \xr{dynamophagy}.

\E{organourgy}{ˈɔːɡənəʊɜːdʒi}{n.}{Gk. ὄργανον tool + -ουργία working} The making of tools; the building of instruments to build other things. --- \dv{organourge} \textit{n.}

\E{ostracophobia}{ˌɒstrəkəˈfəʊbiə}{n.}{Gk. ὀστρακισμός ostracism + φόβος fear} The fear of being ostracised on the Agora; the silent cause of a thousand unspoken opinions. --- \dv{ostracophobe} \textit{n.}

\E{ouranarch}{ˈʊərənɑːk}{n.}{Gk. οὐρανός heaven + ἀρχός ruler} A ruler of the skies; a master of launch and orbit. --- \dv{ouranarchy} \textit{n.} Cf. \xr{jaculatrix}.

\E{oxyrhythmia}{ˌɒksiˈrɪðmiə}{n.}{Gk. ὀξύς sharp + ῥυθμός rhythm} A rapid tempo of releases; the accelerated heartbeat of a lab shipping weekly. --- \dv{oxyrhythmic} \textit{adj.}
